% The clean side of the `linkage paper` fixture pair: a whole small paper that every
% check passes. Its twin, ../seeded/, is this paper with one defect seeded per check,
% each marked `% SEED: <tag>`. Keeping the two in step is the point — a change here that
% makes the clean paper fail is as much a finding as a seeded defect going unreported.
\documentclass[11pt]{article}
\usepackage{amsmath,amsthm}

\newtheorem{theorem}{Theorem}[section]
\newtheorem{proposition}[theorem]{Proposition}
\newtheorem{lemma}[theorem]{Lemma}

\title{A Small Paper About Kernels}
\author{A. Author}

\begin{document}
\maketitle

\begin{abstract}
We study a family of smoothing kernels on the half-line and characterise the members
that respect a prescribed scale covariance. The setting is elementary: a kernel is a
measurable map, the family is indexed by a scale parameter, and the composition of two
stages of smoothing is required to be a third stage of the same family. Under that
requirement alone we show that the admissible kernels form a one-parameter class, that
each member is determined by its first moment, and that the class is closed under the
limits one would want it to be closed under. Two consequences are recorded. The first
is a bound on the tail of every admissible kernel, uniform in the scale parameter and
sharp at the endpoint. The second is a criterion, checkable on the first two moments
alone, for a candidate kernel to belong to the class at all. Neither consequence needs
the smoothness that earlier treatments assumed, and both are stated so that a reader who
cares only about the criterion may stop after the first section of the development. We
close with the two questions the elementary argument leaves open, and with a worked
example that shows the bound is attained.
\end{abstract}

\section{The kernels}
\label{sec:kernels}

A kernel is a measurable map from the half-line to itself, following
\cite{author2020kernels}. The scale parameter is written $t$, and the family is
admissible when composition stays inside it.

\begin{proposition}
  \label{prop:two}
  Let $K_t$ be admissible. Then
  \begin{enumerate}
    \item $K_t$ has a finite first moment, and
    \item $K_t$ is determined by that moment.
  \end{enumerate}
\end{proposition}

By Proposition~\ref{prop:two}(2) the family is a curve rather than a surface, which is
what makes the classification of Section~\ref{sec:consequences} finite.

\begin{lemma}
  \label{lem:moment}
  The first moment of $K_t$ is affine in $t$.
\end{lemma}

\section{Consequences}
\label{sec:consequences}

Lemma~\ref{lem:moment} and Proposition~\ref{prop:two}(1) together give the tail bound,
which is the form \cite{writer2019tails} states for the stable case:
\begin{equation}
  \int_x^\infty K_t(s)\,\mathrm{d}s \le C\,\mathrm{e}^{-x/t}.
  \tag{2.1}
\end{equation}

\begin{theorem}
  \label{thm:criterion}
  A candidate kernel is admissible if and only if its first two moments satisfy the
  relation of Lemma~\ref{lem:moment}.
\end{theorem}

Theorem~\ref{thm:criterion} is the criterion promised in the abstract; the worked
example of \cite{author2020kernels} attains the bound.


\section*{Author contributions}
The author wrote the paper.

\section*{Competing interests}
The author declares none.

\section*{Data availability}
No data were generated.

\section*{Funding}
This work received no funding.

\bibliographystyle{plain}
\bibliography{refs}

\end{document}
